%% kapa: exam=1 round=19 year=2008 period=1 form=A subject=accounting subject_ko=회계학 seq=37 number=117
%% source: Q-Net 감정평가사 기출문제 (cst00309) artlSeq=1001962 file=extracted/1008/1001962/2008년 제19회 감정평가사 1차시험 기출문제/1교시.hwp
%% conversion: kordoc 4.24.51 -> tools/gen1.py (auto); review=proofread-sonnet; flags=none
\begin{question}{117}
(주)태양은 20X7년 1월 1일(리스 실행일) (주)동방리스로부터 다음의 조건으로 기계 1대를 리스하였다.
\begin{framedtext}
․ 리스한 기계의 공정가치는 ￦420,000이고, 내재이자율은 연 10\%이다. \par
․ 리스한 기계의 최소리스료의 현재가치는 공정가치와 일치한다. \par
․ 리스기간은 5년이고, 리스기간 종료시 리스한 기계의 보증잔존가치는 없는 것으로 한다. \par
․ 리스한 기계의 내용연수는 6년이고, 잔존가치는 없으며 감가상각은 정액법으로 한다.
\end{framedtext}
(주)태양이 이상의 리스계약에 대하여 적절한 회계처리를 한다면, 이 리스계약에 따라 20X7년에 인식할 비용은 얼마인가? (단, 법인세효과는 무시한다.) \par
\begin{choices}
  \item ￦42,000
  \item ￦70,000
  \item ￦84,000
  \item ￦112,000
  \item ￦126,000
\end{choices}
\end{question}
\answer{5}
